\section{Q2：如何让 AI 提出有用的研究问题？}
\label{sec:questions}

研究问题说明什么值得弄清，研究提案还说明准备如何弄清。很多系统由人类提供问题或动机，再在这一背景中生成新方法；输出可以有用，但不代表自主发现问题。本文认为，有用的问题应说明尚未解决的主张、区分合理的替代解释，并指出哪些观察会改变答案。这一标准也为有信息量的负结果保留空间，不将科研质量等同于提出能提高 benchmark 分数的方法。

\subsection{候选问题从哪里来}

基于文献的方法主要区别在于如何组织搜索空间。SciMON 从语义、知识图谱和引用邻域获取启发，再对照相关文献修改想法；其输入背景已经给出动机和约束 \cite{scimon}。ResearchAgent 围绕种子论文及关联知识发展问题、方法和实验计划；SciPIP 结合检索与 brainstorming；Chain of Ideas 则把已有发展组织成研究轨迹 \cite{researchagent,scipip,chain-of-ideas}。这些方法为已有研究的可能延伸设置了不同先验。更结构化的检索可以提高相关性，也可能把注意力限制在文献充分的方向。

扩大输入范围并不能消除对表示方式的依赖。MOOSE 从网页语料提取背景与启发，生成社会科学假设；SciAgents 通过本体图谱与专家 agents 连接概念 \cite{moose,sciagents}。它们的机制与通用构思有关，但对社会科学假设的判断、仿生材料领域的展示，并不能证明向 AI 研究迁移。尤其是，图中的关联只提供候选类比，不构成所提机制有效的证据。

数据驱动方法引入了目标更容易观察的反馈循环。HypoGeniC 维护自然语言假设，用带标签样本检验预测，再利用错误和受 bandit 启发的规则更新；HypoRefine 将文献假设与这一过程结合，比较假设修正及集合合并等方式 \cite{hypogenic,hyporefine}。分类场景允许留出集预测评价，为语言主张与观察结果建立连接，但预测规律可能仍只是相关性，而且研究问题由外部给定，并非系统自主发现。

\subsection{人类控制与重组单位}

Scideator 将论文拆为目的、机制和评价 facets，让用户选择、检查新颖性证据并修改组合。其被试内研究在相同底层 LLM 下比较 facet 表示与论文级交互 \cite{scideator}，因此证据对应界面和表示方式，而不是笼统归因于更强模型。这也为人类保留了判断陌生类比及其相关性的空间。

IdeaSynth 同样通过文献反馈支持 facet 的变化与组合，CoQuest 则在人机互动中考察 breadth-first 与 depth-first 的问题发展方式 \cite{ideasynth,coquest}。它们提示，构思常常是反复调整约束的过程，并非一次输出完整问题。不过，用户控制、感知到的创造力支持和最终科学效用是不同结果。帮助研究者探索更多方向的工具，并不一定学会了选择值得执行的方向。

\begin{table}[t]
\centering
\small
\caption{不同构思场景暴露了问题形成过程的不同部分。}
\label{tab:question-settings}
\begin{tabularx}{\linewidth}{@{}P{0.23\linewidth}YY@{}}
\toprule
场景 & 机制或证据 & 尚未解决的关键步骤 \\
\midrule
文献延伸 \cite{scimon,researchagent,scipip,chain-of-ideas} & 检索并重新组织已有研究 & 判断新延伸是否重要 \\
跨概念假设 \cite{moose,sciagents} & 提取启发或遍历图谱 & 验证所提出的关系 \\
数据驱动假设 \cite{hypogenic,hyporefine} & 依据预测反馈更新主张 & 从相关性走向有信息量的问题 \\
人机共同构思 \cite{scideator,ideasynth,coquest} & 用户操作问题组成部分 & 区分人类选择与系统能力 \\
想法 benchmark \cite{ideabench,liveideabench} & 参考文献背景或极少关键词 & 将感知创造力连接到可执行探究 \\
\bottomrule
\end{tabularx}
\end{table}

\subsection{生成质量与筛选质量是两个问题}

NLP 构思的大规模人类研究发现，LLM 提案可以获得较好的新颖性判断，但可行性较弱；进一步分析还发现，扩大生成规模会出现重复想法，模型排序也不可靠 \cite{llm-novel-research-ideas-human-study}。这一发现意味着，大候选池只有在增加了有用替代方案、且 selector 能识别它们时才有帮助。流畅文本或更高采样温度都不能保证这两个条件。

IdeaBench 在参考论文背景下通过模型比较评价生成，LiveIdeaBench 则用极少关键词测试 divergent thinking \cite{ideabench,liveideabench}。由于信息条件与目标行为不同，两者的分数回答不同问题。探索时发散很有价值，收敛到有依据、可检验的问题则是另一项要求。因此，比较应控制候选数量、信息访问和筛选预算，并分别报告候选池多样性与入选提案质量。

\subsection{局限与开放问题}

核心缺口在于：合理的提案不一定对应能解决未知的实验。文献新颖性检查、用户研究和预测测试各覆盖一部分，但没有一种能同时建立重要性、因果有效性和可行性。如何不继承既有 framing 而发现问题，如何超越文本 embedding 识别有意义的多样性，以及何时放弃吸引人但不可检验的想法，仍待研究。时间受限的文献检索也需要配套的模型污染审计。以下均为拟议研究，其预期优势仍是假设。

\subsection{拟议计划一：同时生成提案与有区分力的可执行 pilot}

\paragraph{研究假设与方法。}
要求提案包含替代解释及可执行 pilot，有望提高真正解决某个明确未知的项目比例。生成器必须给出核心主张、合理替代解释、两种解释各自预期的观察，以及能够区分它们的最小比较。例如，声称校准改善时，应区分方法机制与额外调参的影响。固定 coding agent 实现 pilot，另设 checker 检查测量是否对应所声明的解释。

\paragraph{任务与划分。}
使用十二个有边界的 AI 研究情境，覆盖冻结小模型的校准、检索评分和受控输入变化下的鲁棒性。六个用于开发提示和检查实验基础设施，另六个采用不相交的数据集与代码模板家族并封存。允许开发集 pilot 反馈，但在独立评价种子和测试样本开放前冻结主张与分析规则；pilot 不得访问最终测试集。

\paragraph{基线、指标与风险。}
比较普通提案、带文字实验计划的提案，以及带有区分力的可执行 pilot 的提案；匹配总推理与实验预算，pilot 执行计入后者。分别消融替代解释和独立 checker。测量可执行性、盲评的区分能力、可复现的主张判定、置信度校准和每个有信息量结果的成本，将无法定论与实现错误分开。风险包括偏爱容易问题、pilot 统计效力不足及 checker 奖励格式合规；脚本成功运行但比较无效，不能算成功。

\subsection{拟议计划二：将探索能力与筛选能力分开测量}

\paragraph{研究假设与方法。}
机制层面的多样性只有与能识别有用差异的 selector 配合，才会改善下游产出。将普通采样、检索生成和受控 facet 重组三种生成器，与随机、直接提示、证据驱动三类 selector 交叉比较。每种生成器在相同 token 与检索预算内产出二十四个候选，冻结候选池后再应用各 selector，以区分生成、筛选及二者交互的影响，而非只比较完整 pipeline。

\paragraph{数据、消融与测量。}
复用上述开发与封存情境，在评价留出候选池前固定所有选择。盲态标注者依据实质机制聚类并识别有效替代方案，embedding 仅辅助标注。对每个 selector 的前三个候选及独立随机审计样本执行相同 pilot 协议。比较仅扩大候选数量、embedding 去重，以及去掉新颖性检查的 facet 重组。报告有效机制覆盖、重复负担、top-$k$ 有信息量产出和成本；记录抽样概率的随机审计有助于发现所有 selector 都遗漏的候选。

\paragraph{失败风险。}
机制标签可能主观，selector 也可能偏爱自身生成器的语言风格。应交叉使用不同模型，报告标注分歧，并在保留原始表达之外增加规范化评价条件。由于仅执行候选子集，不能声称得知候选池的 oracle 最优解或精确 selection regret。若文本聚类数量增加却未提高经验证的产出，则不支持研究假设。
